\subsection{Basis of the Clifford algebra.}

PROPOSITION 3. --- Let Q and Q' be two quadratic forms and F a bilinear form on E such that Q'(x) = Q(x) + F(x, x) for all x ∈ E. The map λ\_F maps the ideal I(Q') onto the ideal I(Q), and defines an isomorphism (denoted \(\overline{\lambda_{\mathbf{F}}}\) ) from the A-module C(Q') onto the A-module C(Q).

As \(\lambda_{\mathbb{F}}\) is a bijection whose \(\lambda_{-\mathbb{F}}\) is the inverse bijection (Lemma 3), it suffices to prove the inclusion \(\lambda_{\mathbb{F}}(\mathrm{I}(\mathrm{Q}^{\prime})) \subset \mathrm{I}(\mathrm{Q})\) . Since \(\mathrm{I}(\mathrm{Q})\) is a left ideal stable under \(i_x^{\mathbb{F}}\) (Lemma 1), (8) shows that the set of \(u \in \mathrm{T}(\mathrm{E})\) such that \(\lambda_{\mathbb{F}}(u) \in \mathrm{I}(\mathrm{Q})\) is a left ideal. It therefore suffices to show that, for any \(u \in \mathrm{T}(\mathrm{E})\) and \(x \in \mathrm{E}\) , one has \(\lambda_{\mathbb{F}}(x \otimes x \otimes u - \mathrm{Q}'(x)u) \in \mathrm{I}(\mathrm{Q})\) . Now, by (8) and Lemma 1, we have

\[
\lambda_ {\mathrm{F}} \circ e _ {x} ^ {2} = (e _ {x} + i _ {x} ^ {\mathrm{F}}) ^ {2} \circ \lambda_ {\mathrm{F}} = (e _ {x} ^ {2} + \mathrm{F} (x, x)) \circ \lambda_ {\mathrm{F}},
\]

whence

\[
\begin{array}{r l} \lambda_ {\mathbf {F}} (x \otimes x \otimes u - \mathrm{Q} ^ {\prime} (x) u) & = (e _ {x} ^ {2} + \mathrm{F} (x, x) - \mathrm{Q} ^ {\prime} (x)) \circ \lambda_ {\mathbf {F}} (u) \\ & = (x \otimes x - \mathrm{Q} (x)) \otimes \lambda_ {\mathbf {F}} (u) \in \mathrm{I} (\mathrm{Q}). \end{array}
\]

Lemma 4. — If the quadratic form Q is zero, then C(Q) is none other than the exterior algebra of E.

Indeed, the exterior algebra of E is none other than the quotient of T(E) by the two-sided ideal J generated by the elements of the form \(a \otimes v \otimes a\) ( \(a \in E\) , \(v \in T(E)\) ) (chap.~III, § 5, no.~5 and no.~9). It is clear that I(Q) ⊂ J. It therefore suffices to show that \(a \otimes v \otimes a \in I(Q)\) . This is evident if \(v \in T^0\) . Suppose this assertion has been proved for \(n-1 \leqslant \nu \in \sum_{h=0}^{\infty} \mathrm{T}^{h}\) , and let \(x \in \mathrm{E}\) and \(u \in \mathrm{T}^{n-1}\) ; we have

\[
\begin{array}{c} \boldsymbol {a} \otimes \boldsymbol {x} \otimes \boldsymbol {u} \otimes \boldsymbol {a} = (\boldsymbol {a} + \boldsymbol {x}) \otimes (\boldsymbol {a} + \boldsymbol {x}) \otimes \boldsymbol {u} \otimes \boldsymbol {a} \\ - \boldsymbol {a} \otimes \boldsymbol {a} \otimes \boldsymbol {u} \otimes \boldsymbol {a} - \boldsymbol {x} \otimes \boldsymbol {a} \otimes \boldsymbol {u} \otimes \boldsymbol {a} - \boldsymbol {x} \otimes \boldsymbol {x} \otimes \boldsymbol {u} \otimes \boldsymbol {a} \end{array}
\]

and the four terms on the right-hand side belong to I(Q).

Suppose in particular that the module E admits a basis \((x_{i})_{i\in \mathbf{L}}\) , and let us totally order the index set L. We know (chap.~III, §5, no.~6) that the exterior algebra of E admits as a basis the family consisting of the elements \(x_{\mathrm{H}}\) where H ranges over the set of finite subsets of L and where \(x_{\mathrm{H}}\) is the element \(x_{h_1}\wedge \ldots \wedge x_{h_q}\) , \((h_1,\ldots ,h_q)\) denoting the strictly increasing sequence of elements of H.

On the other hand, consider the bilinear form F defined by \(\mathrm{F}(x_{i}, x_{j}) = -\Phi(x_{i}, x_{j})\) if i \textgreater{} j, \(\mathrm{F}(x_{i}, x_{j}) = 0\) if i \textless{} j and \(\mathrm{F}(x_{i}, x_{i}) = -\mathrm{Q}(x_{i})\) . It is clear that \(\mathrm{Q}(x) + \mathrm{F}(x, x) = 0\) ; with the notations of Prop. 3, it follows from this proposition and Lemma 4 that \(\bar{\lambda}_{\mathrm{F}}\) is an isomorphism of \(\wedge \mathrm{E}\) onto C(Q), which is therefore a free A-module. Let us prove, by induction on the number of elements of H, that we have

\[
\overline {{{{\lambda}}}} _ {\mathrm{F}} (x _ {\mathrm{H}}) = \rho (x _ {h _ {1}}) \dots \rho (x _ {h _ {q}})\tag{11}
\]

\((\mathrm{where}\ (h_{1},\ldots,h_{q})\) is the strictly increasing sequence of elements of H). This is evident if H has 0 or 1 element. Suppose (11) holds for sets having at most \(q-1\) elements. Consider a subset H of \(q\) elements, denote \(j\) its smallest element, and let us write \(\mathrm{H}=\{j\}\cup\mathrm{K}\) where K is a subset with \(q-1\) elements. By (8) and the induction hypothesis, we have

\[
\overline {{{\lambda}}} _ {\mathrm{F}} (x _ {\mathrm{H}}) = \overline {{{\lambda}}} _ {\mathrm{F}} (x _ {j} \wedge x _ {\mathrm{K}}) = \rho (x _ {j}) \overline {{{\lambda}}} _ {\mathrm{F}} (x _ {\mathrm{K}}) + i _ {x _ {j}} ^ {\mathrm{F}} (\overline {{{\lambda}}} _ {\mathrm{F}} (x _ {\mathrm{K}})) = x _ {\mathrm{H}} ^ {\prime} + i _ {x _ {j}} ^ {\mathrm{F}} (x _ {\mathrm{K}} ^ {\prime}),
\]

setting, for every finite subset J of L, \(x_{J}^{\prime}=\rho(x_{j_{1}})\ldots\rho(x_{j_{s}})\) , where \((j_{1},\ldots,j_{s})\) is the increasing sequence of elements of J. Now, for \(i\in\mathbf{K}\) , \(x_{i}\) belongs to the kernel of the linear form \(y\to\mathrm{F}(x_{j},y)\) , hence \(i_{x_{j}}^{\mathrm{F}}(x_{\mathrm{K}}^{\prime})=0\) (Lemma 1). This proves the desired result.

We have therefore proved the following theorem:

THEOREM 1. --- Suppose that the A-module E admits a basis \((x_{i})_{i\in L}\) , the index set L being equipped with a total order structure. For every finite subset H of L, set \(x_{\mathrm{H}} = \rho(x_{h_1})\rho(x_{h_2})\ldots \rho(x_{h_q})\) , where \((h_1,\ldots,h_q)\) is the strictly increasing sequence of the elements of H. Then the elements \(x_{\mathrm{H}}\) form a basis of the A-module C(Q).

COROLLARY 1. --- If E is a free module of dimension n, then C(Q) is a free module of dimension \(2^{n}\) ; moreover, if \(n > 0\) , C \(^{+}\) and C \(^{-}\) are free modules of dimension \(2^{n-1}\) .

This follows immediately from the properties of binomial coefficients.

COROLLARY 2. --- If E is a free module, the canonical mapping \(\rho\) from E into C(Q) and the mapping \(a \to a\) .1 of A into C(Q) are injective.

COROLLARY 3. --- Suppose that E is the direct sum of two free submodules \(E_{1}\) and \(E_{2}\) . Let \(Q_{i}\) be the restriction of Q to \(E_{i}\) and \(p_{i}\) the canonical map from C(Q \(_{i}\) ) into C(Q) (i = 1, 2). Then the

linear map \(p\) of \(\mathrm{C(Q_{1})\otimes C(Q_{2})}\) into \(\mathrm{C(Q)}\) deduced from the bilinear map \((a,b)\to p_{1}(a)p_{2}(b)\) of \(\mathrm{C(Q_{1})\times C(Q_{2})}\) in \(\mathrm{C(Q)}\) is a bijection.

It suffices indeed to consider the basis of E obtained by taking the union of a basis of \(E_{1}\) and a basis of \(E_{2}\) .

COROLLARY 4. --- With the hypotheses and notation of Cor. 3, suppose further that \(E_{1}\) and \(E_{2}\) are orthogonal, and transport to \(C(Q_{1})\otimes C(Q_{2})\) , by means of the bijection p, the algebra structure of \(C(Q)\) . If \(a_{i}\) and \(b_{i}\) are even or odd elements of \(C(Q_{i})\) (i=1, 2), we have \((a_{1}\otimes a_{2})(b_{1}\otimes b_{2})=\varepsilon(a_{1}b_{1})\otimes(a_{2}b_{2})\) , with \(\varepsilon=1\) unless \(a_{2}\) and \(b_{1}\) are odd, in which case \(\varepsilon=-1\) .

It suffices indeed to prove that \(p_2(a_2)p_1(b_1) = \varepsilon p_1(b_1)p_2(a_2)\) , and for that one may assume that \(p_2(a_2)\) (resp. \(p_1(b_1)\) ) is a product \(x_1 \ldots x_h\) (resp. \(y_1 \ldots y_k\) ) of elements of \(\rho_{\mathbb{Q}}(\mathbf{E}_2)\) (resp. \(\rho_{\mathbb{Q}}(\mathbf{E}_1)\) ). Since \(\mathbf{E}_1\) and \(\mathbf{E}_2\) are orthogonal, one has

\[
x _ {i} y _ {j} + y _ {j} x _ {i} = \Phi (x _ {i}, y _ {j}) = 0,
\]

whence

\[
x _ {1} \dots x _ {h} y _ {1} \dots y _ {k} = (- 1) ^ {h k} y _ {1} \dots y _ {k} x _ {1} \dots x _ {h}.
\]

The conclusions of Cor. 3 and 4 remain true if one omits the hypothesis that \(\mathbf{E}_1\) and \(\mathbf{E}_2\) are free modules (cf.~exerc. 1).

